\lstset{language = matlab}
\lstset{numbers = left, numberstyle = \tiny, stepnumber = 2, numbersep = 5pt}
\lstset{tabsize = 3}
%\lstset{index={Condition!normal}}

% Main Program
\begin{lstlisting}[caption = {Main Program}, label = lst:Main_Program]

clear;

% Export the figure to the following file.
% File_Name = 'Untitled';

% Normal Condition on the left hand side.
% P = 1;          % for q(x) equal to poly1 and poly 2
P = 0.01;       % for an exponential q(x)

% Dirichlet condition on theleft hand side.
delta = 0;

% Degree of denominator in q(x):
% q_Denom_Type = 'Poly_1';
% q_Denom_Type = 'Poly_2';
q_Denom_Type = 'Exp';

% Coefficients of the load function:
% load_Coeff = 1;
% load_Coeff = [1 2];
load_Coeff = [1 2 -3];

% Coefficients of q:
% q_Coeff = [10 0.01]; % For q(x) with linear and quadratic denominator
q_Coeff = 1;           % For the exponential choice

% No. of elements:
N = [4 8 16];

[h, elem_Size, sol_Size, q_Func, load_Func, ...
		x, u_FEM_Lin, u_FEM_Cub, u_Exact] = ...
		Def_Problem( N, q_Denom_Type, load_Coeff, q_Coeff, delta, P );

% Calculate FEM solution:
for i = 1:1:length(N)
		[u_FEM_Lin(:, i), u_FEM_Cub(:, i)] = ...
				Solve_FEM_Problem( N( i ), h( i ), ...
				elem_Size( i ), sol_Size, x, P, delta, q_Func, load_Func );
end;

% Error estimation:
[conv_Factor, sq_Error] = ...
		Estimate_Error( N, sol_Size, x, u_Exact, u_FEM_Lin, u_FEM_Cub );

\end{lstlisting}
\line(1,0){450}
\vspace{5 mm}

% Define the problem and its analytical solution:
\begin{lstlisting}[caption = {Define the Problem and Its Analytical Solution}, label = lst:Def_Problem]

function [h, elem_Size, sol_Size, q_Function, load_Func, ...
		x, u_FEM_Lin, u_FEM_Cub, u_Exact] = ...
		Def_Problem( N, q_Denom_Type, load_Coeff, q_Coeff, delta, P )

% A uniform discretizatin
h = 1 ./ N;

% No. of points in the final plot:
sol_Size = 2^16;

% Element size in the FEM solution
elem_Size = sol_Size ./ N;

% Evaluation points in the final solution
x = linspace(0, 1, sol_Size + 1 )';

% Initialize FEM-solution
u_FEM_Lin = zeros( sol_Size + 1, length(N) );
u_FEM_Cub = zeros( sol_Size + 1, length(N) );
load_Degree = length(load_Coeff) - 1;

% Convert coefficients of the polynomial load to a general
% quadratic polynomial.
if load_Degree == 0
		load_Coeff = [load_Coeff 0 0];
elseif load_Degree == 1
		load_Coeff = [load_Coeff 0];
end;

if strcmpi(q_Denom_Type, 'Poly_1') == 1
		[q_Function, load_Func, u_Exact] = ...
		Def_q_Denom_1( load_Coeff, q_Coeff, delta, P );

elseif strcmpi(q_Denom_Type, 'Poly_2') == 1
		[q_Function, load_Func, u_Exact] = ...
		Def_q_Denom_2( load_Coeff, q_Coeff, delta, P );

elseif strcmpi(q_Denom_Type, 'Exp') == 1
		[q_Function, load_Func, u_Exact] = ...
		Def_q_Exp( load_Coeff, q_Coeff, delta, P );
end;


function [q_Function, load_Func, u_Exact] = ...
		Def_q_Denom_1( load_Coeff, q_Coeff, delta, P )

d = zeros(1, 6);
q_Function = @(x) q_Coeff(1) ./ ( x + q_Coeff(2) );
load_Func = @(x) load_Coeff(1) + load_Coeff(2) * x ...
		+ load_Coeff(3) * x.^2;

d(1) = delta;

d(2) = q_Coeff(2) * P / q_Coeff(1) ...
			 + q_Coeff(2) * load_Coeff(1) / q_Coeff(1) ...
			 + q_Coeff(2) * load_Coeff(2) / ( 2 * q_Coeff(1) ) ...
			 + q_Coeff(2) * load_Coeff(3) / ( 3 * q_Coeff(1) );

d(3) = P / ( 2 * q_Coeff(1) ) ...
		+ ( 1 - q_Coeff(2) ) * load_Coeff(1) / ( 2 * q_Coeff(1) ) ...
		+ load_Coeff(2) / ( 4 * q_Coeff(1) ) ...
		+ load_Coeff(3) / ( 6 * q_Coeff(1) );        

d(4) = - load_Coeff(1) / ( 3 * q_Coeff(1) ) ...
		- q_Coeff(2) * load_Coeff(2) / ( 6 * q_Coeff(1) );

d(5) = - load_Coeff(2) / ( 8 * q_Coeff(1) ) ...
		- q_Coeff(2) * load_Coeff(3) / ( 12 * q_Coeff(1) );

d(6) = -load_Coeff(3) / ( 15 * q_Coeff(1) );

u_Exact = @(x) d(1) + d(2)*x + d(3)*x.^2 + d(4)*x.^3 + ...
		d(5)*x.^4 + d(6)*x.^5;


function [q_Function, load_Func, u_Exact] = ...
		Def_q_Denom_2( load_Coeff, q_Coeff, delta, P )
		
d = zeros(1, 7);
q_Function = @(x) q_Coeff(1) ./ ( x.^2 + q_Coeff(2) );    
load_Func = @(x) load_Coeff(1) + load_Coeff(2) * x ...
		+ load_Coeff(3) * x.^2;

d(1) = delta;

d(2) = q_Coeff(2) * P / q_Coeff(1) + ...
		q_Coeff(2) * load_Coeff(1) / q_Coeff(1) + ...
		q_Coeff(2) * load_Coeff(2) / ( 2 * q_Coeff(1) ) + ...
		q_Coeff(2) * load_Coeff(3) / ( 3 * q_Coeff(1) );

d(3) = -q_Coeff(2) * load_Coeff(1) / ( 2 * q_Coeff(1) );

d(4) = P / ( 3 * q_Coeff(1) ) ...
		+ load_Coeff(1) / ( 3 * q_Coeff(1) ) ...
		+ ( 1 - q_Coeff(2) ) * load_Coeff(2) / ( 6 * q_Coeff(1) ) ...
		+ load_Coeff(3) / ( 9 * q_Coeff(1) );

d(5) = -load_Coeff(1) / ( 4 * q_Coeff(1) ) ...
		- q_Coeff(2) * load_Coeff(3) / ( 12 * q_Coeff(1) );

d(6) = - load_Coeff(2) / ( 10 * q_Coeff(1) ); 

d(7) = - load_Coeff(3) / ( 18 * q_Coeff(1) );

u_Exact = @(x) d(1) + d(2)*x + d(3)*x.^2 + d(4)*x.^3 ...
		+ d(5)*x.^4 + d(6)*x.^5 + d(7)*x.^6;


function [q_Function, load_Func, u_Exact] = ...
				Def_q_Exp( load_Coeff, q_Coeff, delta, P )

% An exponential q(x) and a general up to 2nd degree load polynomial:
d = zeros(1, 5);
alpha = q_Coeff(1);
q_Function = @(x) exp( -alpha * x );
load_Func = @(x) load_Coeff(1) + load_Coeff(2)*x + load_Coeff(3).*x.^2;

d(1) = delta ...
		- load_Coeff(1) / alpha ...
		- load_Coeff(2) / (2 * alpha) ...
		- load_Coeff(3) / (3 * alpha) ...
		- load_Coeff(1) / alpha^2 ...
		+ load_Coeff(2) / alpha^3 ...
		- 2 * load_Coeff(3) / alpha^4 ...
		- P / alpha;

d(2) = load_Coeff(1) / alpha ...
		+ load_Coeff(2) / ( 2 * alpha )...
		+ load_Coeff(3) / (3 * alpha) ...
		+ load_Coeff(1) / alpha^2 ...
		- load_Coeff(2) / alpha^3 ...
		+ 2 * load_Coeff(3) / alpha^4 + P/alpha;

d(3) = - load_Coeff(1) / alpha ...
		+ load_Coeff(2) / alpha^2 ...
		- 2 * load_Coeff(3) / alpha^3;

d(4) = - load_Coeff(2) / ( 2 * alpha ) ...
		+ load_Coeff(3) / alpha^2;

d(5) = - load_Coeff(3) / ( 3 * alpha );

u_Exact = @(x) d( 1 ) ...
		+ ( d(2) + d(3)*x + d(4)*x.^2 + d(5)*x.^3 ) .* exp(alpha * x);

\end{lstlisting}
\line(1,0){450}
\vspace{5 mm}

% Set up and solve Equation System. Then buildup the solution
\begin{lstlisting}[caption = {Solve FEM-Problem and Build up Solution as a Power Series}, label = lst:Solve_Eq_Sys]

function [u_FEM_Lin, u_FEM_Cub] = Solve_FEM_Problem( N, h, elem_Size, ...
		sol_Size, x, P, delta, q_Func, load_Func )

% Define local and global basis functions and their first derivative.
[theta_Lin, theta_Prime_Lin, psi_Lin, psi_Prime_Lin, ...
		theta_Cub, theta_Prime_Cub, psi_Cub, psi_Prime_Cub] = ...
		Def_FEM_Func( delta );

% Set up the Eq. system
[K_Lin, b_Lin, K_Cub, b_Cub] = Setup_Eq_Sys( N, h, P, q_Func, load_Func, ...
		theta_Prime_Lin, psi_Lin, psi_Prime_Lin, ...
		theta_Prime_Cub, psi_Cub, psi_Prime_Cub );

% Solve both systems.
a_Lin = K_Lin\b_Lin;
a_Cub = K_Cub\b_Cub;

% Construct FEM-Solution as a linear combination of basis-functions.
[u_FEM_Lin, u_FEM_Cub] = Build_FEM_Sol( N, elem_Size, x, ...
		sol_Size, a_Lin, a_Cub, theta_Lin, theta_Cub, psi_Lin, psi_Cub );

\end{lstlisting}
\line(1,0){450}
\vspace{5 mm}


% Estimate the numerical error
\begin{lstlisting}[caption = {Solution Visualization, Error Estimation and Eventual Convergence}, label = lst:Estimate_Error]

function [max_Error_Lin, sq_Error_Lin, max_Error_Cub, sq_Error_Cub] = ...
    Estimate_Error( N, sol_Size, x, u_Exact, u_FEM_Lin, u_FEM_Cub, ...
    varargin )

u_Exact = u_Exact(x);

% Plot both Exact and Finite Element solutions
n = length( N );
u_FEM = zeros( sol_Size + 1, n, 2 );
u_FEM(:, :, 1) = u_FEM_Lin;
u_FEM(:, :, 2) = u_FEM_Cub;

max_Error = zeros(2, n);
sq_Error = zeros(2, n);

h = figure();
left = 0.09;
bottom = 0.55;
width = 0.24;
height = 0.40;
horizontalSpace = 0.08;

for k = 1:1:2
    for m = 1:1:n
        subplot(2, n, (k - 1)*n + m, 'Position', ...
            [left bottom width height] );
        
        plot(x, u_Exact , 'red', x, u_FEM(:, m, k), ...
            'green', 'LineWidth', 2);
        
        left = left + width + horizontalSpace;

        % Grid and thick marks:
        grid;
        set( gca, 'FontName', 'Arial', 'FontSize', 14 );
        set( gca, 'XTick', [0.25 0.5 0.75 1] );
                
        % Legend
        str_FEM = ['$u_{cub}$, ' 'N = ' num2str( N(m) ) ];
        if k == 1
            str_FEM = ['$u_{lin}$, ' 'N = ' num2str( N(m) ) ];
        end;

        h_leg = legend( '$u_{ex}$', str_FEM, ...
            'Location', 'Northwest', 'Orientation', 'Vertical' );
        set( h_leg, 'Interpreter', 'Latex', 'FontSize', 14 );

        abs_Error = abs( u_Exact - u_FEM( :, m, k ) );
        max_Error(k, m) = max( abs_Error );
        sq_Error(k, m) = sqrt( sum( abs_Error.^2 ) );

    end;
    
    left = 0.09;
    bottom = 0.05;
    width = 0.24;
    height = 0.40;
    horizontalSpace = 0.08;
end;

% Linear and cubic errors compared to exact solution:
max_Error_Lin = max_Error(1, :);
sq_Error_Lin = sq_Error(1, :);
max_Error_Cub = max_Error(2, :);
sq_Error_Cub = sq_Error(2, :);

% Print figure
if nargin == 7
    paperWidth = 21;
    paperHeight = 29.7;
    Export_Figure( h, paperWidth, paperHeight, varargin{end} );
end;

\end{lstlisting}
\clearpage

% Define Finite Element Basis-Functions
\begin{lstlisting}[caption = {Linear and Cubic Shape-Functions}, label = lst:Def_FEM_Func]

function [theta_Lin, theta_Prime_Lin, psi_Lin, psi_Prime_Lin, ...
    theta_Cub, theta_Prime_Cub, psi_Cub, psi_Prime_Cub] = ...
    Def_FEM_Func( delta )

[theta_Lin, theta_Prime_Lin, psi_Lin, psi_Prime_Lin] = ...
        Def_Linear_FEM(delta);

[theta_Cub, theta_Prime_Cub, psi_Cub, psi_Prime_Cub] = ...
    Def_Cubic_FEM(delta);

function [theta_Lin, theta_Prime_Lin, psi_Lin, psi_Prime_Lin] = ...
        Def_Linear_FEM(delta)

% Linear shape functions, in local coordinates.
psi_Lin = cell( 2, 1 );
psi_Lin{1} = @(y) 1 - y;
psi_Lin{2} = @(y) y;

% Derivative of lineat shape functions in the local coordinates.
psi_Prime_Lin = cell( 2, 1);
psi_Prime_Lin{1} = @(y) -1;
psi_Prime_Lin{2} = @(y) 1;

% The global function theta and its derivative to satisfy Dirichlet
% boundary condition.
theta_Lin = @(y) delta -delta*y;
theta_Prime_Lin = @(y) -delta;

function [theta_Cub, theta_Prime_Cub, psi_Cub, psi_Prime_Cub] = ...
    Def_Cubic_FEM(delta)

% Cubic shape functions, in local coordinates.
psi_Cub = cell(4, 1);
psi_Cub{1} = @(y) 1 - 3*y.^2 + 2*y.^3;
psi_Cub{2} = @(y) 3*y.^2 - 2*y.^3;
psi_Cub{3} = @(y) y -2*y.^2 + y.^3;
psi_Cub{4} = @(y) -y.^2 + y.^3;

% Derivative of cubic shape functions in the local coordinates.
psi_Prime_Cub = cell(4, 1);
psi_Prime_Cub{1} = @(y) -6*y + 6*y.^2;
psi_Prime_Cub{2} = @(y) 6*y -6*y.^2;
psi_Prime_Cub{3} = @(y) 1 - 4*y + 3*y.^2;
psi_Prime_Cub{4} = @(y) -2*y + 3*y.^2;

% The global function theta and its derivative to satisfy Dirichlet
% boundary condition.
theta_Cub = @(y) delta - 3*delta*y.^2 + 2*delta*y.^3;
theta_Prime_Cub = @(y) -6*delta*y + 6*delta*y.^2;

\end{lstlisting}
\line(1,0){450}
\vspace{5 mm}


% Setup Equation System
\begin{lstlisting}[caption = {Setup Equation System}, label = lst:Setup_Eq_Sys]

function [K_Lin, b_Lin, K_Cub, b_Cub] = ...
    Setup_Eq_Sys( N, h, P, q_Func, load_Func, ...
    theta_Prime_Lin, psi_Lin, psi_Prime_Lin, ...
    theta_Prime_Cub, psi_Cub, psi_Prime_Cub )

K_Lin = zeros(N + 1, N + 1);
b_Lin = zeros(N + 1, 1);
K_Cub = zeros(2*N + 2, 2*N + 2);
b_Cub = zeros(2*N + 2, 1);

% Build up the global stiffness matrix and load vector
for n = 1:1:N
    
    % Construct element stiffness matrix and element load vector
    [K_Elem_Lin, b_Elem_Lin, K_Elem_Cub, b_Elem_Cub] = ...
        Elem_Cont( h, n, q_Func, load_Func, ...
        psi_Lin, psi_Prime_Lin, theta_Prime_Lin, ...
        psi_Cub, psi_Prime_Cub, theta_Prime_Cub );

    for i = n:1:n + 1
        for j = n:1:n + 1
            
            % Linear assembly
            K_Lin(i, j) = K_Lin(i, j) + K_Elem_Lin(i - n + 1, j - n + 1);
            
            % Cubic assembly
            K_Cub(i, j) = K_Cub(i, j) + K_Elem_Cub(i - n + 1, j - n + 1);
            K_Cub(i, j + N + 1) = K_Cub(i, j + N + 1) ...
                + K_Elem_Cub(i - n + 1, j - n + 3);
            K_Cub(i + N + 1, j) = K_Cub(i + N + 1, j) ...
                + K_Elem_Cub(i - n + 3, j - n + 1);
            K_Cub(i + N + 1, j + N + 1) = K_Cub(i + N + 1, j + N + 1) ...
                + K_Elem_Cub(i - n + 3, j - n + 3);
            
        end;
	 
        % Linear load assembly
        b_Lin(i) = b_Lin(i) + b_Elem_Lin(i - n + 1);

        % Cubic load assembly
        b_Cub(i) = b_Cub(i) + b_Elem_Cub(i - n + 1);
        b_Cub(i + N + 1) = b_Cub(i + N + 1) + b_Elem_Cub(i - n + 3);
        
    end;    
end;

% Implement boundary conditions:
[K_Lin, b_Lin, K_Cub, b_Cub] = ...
    boundary(N, h, P, q_Func, theta_Prime_Cub, K_Lin, b_Lin, K_Cub, b_Cub);

function [K_Lin, b_Lin, K_Cub, b_Cub] = ...
    boundary(N, h, P, q_Func, theta_Prime_Cub, K_Lin, b_Lin, K_Cub, b_Cub)

% Implementing boundary conditions with linear basis-functions:
% Dirichlet condition at the left boundary:
K_Lin(1, :) = zeros(1, N + 1);
K_Lin(:, 1) = zeros(N + 1, 1);
K_Lin(1, 1) = 1;
b_Lin(1) = 0;

% Normal condition at the right boundary:       
b_Lin(N + 1) = b_Lin(N + 1) + P;

% Implementing boundary conditions with cubic basis-functions:
% Dirichlet condition at the left boundary:
K_Cub(1, :) = zeros(1, 2*N + 2);
K_Cub(:, 1) = zeros(2*N + 2, 1);
K_Cub(1, 1) = 1;
b_Cub(1) = 0;

% Normal condition at the right boundary:
K_Cub(2*N + 2, :) = zeros(1, 2*N + 2);
K_Cub(2*N + 2, 2*N + 2) = 1;
b_Cub(N + 1) = b_Cub(N + 1) + P;
b_Cub(2*N + 2) = h * P / q_Func(1) - h * theta_Prime_Cub(1);

\end{lstlisting}
\line(1,0){450}
\clearpage

% Build Solutions
\begin{lstlisting}[caption = {Build up FEM-Solution}, label = lst:Build_Sol]

function [u_FEM_Lin, u_FEM_Cub] = Build_FEM_Sol( N, elem_Size, x, ...
    sol_Size, a_Lin, a_Cub, theta_Lin, theta_Cub, psi_Lin, psi_Cub )

u_FEM_Lin = Build_Linear_FEM( sol_Size, elem_Size, N, a_Lin, psi_Lin );
u_FEM_Cub = Build_Cubic_FEM( sol_Size, elem_Size, N, a_Cub, psi_Cub );

% Add the global function theta to the solution.
u_FEM_Lin = u_FEM_Lin + theta_Lin(x);
u_FEM_Cub = u_FEM_Cub + theta_Cub(x);


function u_FEM = Build_Linear_FEM(sol_Size, elem_Size, N, a, psi)

% Initialize FEM-Solution:
u_FEM = zeros( sol_Size, 1 );

% Shape-functions in local coordinates:
y = linspace( 0, 1, elem_Size + 1 )';
Gamma = zeros( elem_Size + 1, 2);
for i = 1:1:2
    Gamma(:, i) = psi{i}( y );
end;

% Construct FEM-Solution as a linear combination of the shape-functions.
for e = 1:1:N
    start = (e - 1) * elem_Size + 1;
    last = start + elem_Size;
    u_FEM(start:last) = a(e) * Gamma(:, 1) + a(e + 1) * Gamma(:, 2);
end;


function u_FEM = Build_Cubic_FEM(sol_Size, elem_Size, N, a, psi)

% Initialize FEM-Solution:
u_FEM = zeros( sol_Size, 1 );

% Shape-functions in local coordinates:
y = linspace( 0, 1, elem_Size + 1 )';
Gamma = zeros( elem_Size + 1, 4 );
for i = 1:1:4
    Gamma(:, i) = psi{i}( y );
end;

% Construct a linear combination of the basis functions over all elements.
for e = 1:1:N
    start = (e - 1) * elem_Size + 1;
    last = start + elem_Size;
    u_FEM(start:last) = a(e) * Gamma(:, 1) ...
        + a(e + 1) * Gamma(:, 2) ...
        + a(e + N + 1) * Gamma(:, 3) ...
        + a(e + N + 2) * Gamma(:, 4);
end;

\end{lstlisting}
\line(1,0){450}
\vspace{5 mm}


% Build Solutions
\begin{lstlisting}[caption = {Element Contributions in both Linear and Cubic Case}, label = lst:Elem_Cont]

function [K_Elem_Lin, b_Elem_Lin, K_Elem_Cub, b_Elem_Cub] = ...
        Elem_Cont( h, elem_No, q_Func, load_Func, ...
        psi_Lin, psi_Prime_Lin, theta_Prime_Lin, ...
        psi_Cub, psi_Prime_Cub, theta_Prime_Cub )

% Set up element contribution for the linear basis-functions
[K_Elem_Lin, b_Elem_Lin] = Cal_Cont( 1, h, elem_No, q_Func, ...
    psi_Lin, psi_Prime_Lin, theta_Prime_Lin, load_Func );

% Set up element contribution for the cubic basis-functions
[K_Elem_Cub, b_Elem_Cub] = Cal_Cont( 3, h, elem_No, ...
    q_Func, psi_Cub, psi_Prime_Cub, theta_Prime_Cub, load_Func );


function [K_Elem, b_Elem] = Cal_Cont( basis_Degree, h, elem_No, ...
    q_Func, psi, psi_Prime, theta_Prime, load_Func)

% Initialize linear element contribution for the linear and
% cubic basis-functions.
if basis_Degree == 1
    K_Elem = zeros(2, 2);
    b_Elem = zeros(2, 1);
    stiff_Elem_Int = cell(2, 2);
    load_Int = cell(2, 1);

elseif basis_Degree == 3
    K_Elem = zeros(4, 4);
    b_Elem = zeros(4, 1);
    stiff_Elem_Int = cell(4, 4);
    load_Int = cell(4, 1);
end;

% Calculate entries of element matrix and element load.
for i = 1:1:basis_Degree + 1
    for j = i:1:basis_Degree + 1
        % Define integrands of the element matrix
        stiff_Elem_Int{i, j} = @(y) ...
            ( 1/h ) * ( q_Func( ( elem_No - 1 + y ) * h ) ) ...
            .* ( psi_Prime{i}(y) .* psi_Prime{j}(y) );
        
        % Integrate over the element
        K_Elem(i, j) = quadgk(stiff_Elem_Int{i, j}, 0, 1, ...
            'RelTol', 1e-8, 'AbsTol', 1e-12);
    end;
    
	% Define integrands of the element load
	load_Int{i} = @(y) ...
        h * load_Func( ( elem_No - 1 + y )*h ) .* psi{i}(y) ...
        - h * q_Func( ( elem_No - 1 + y )*h ) .* theta_Prime(y);
    
    % Integrate over the element
    b_Elem(i) = quadgk(load_Int{i}, 0, 1, ...
    'RelTol', 1e-8, 'AbsTol', 1e-12);
end;

% Use symmetrical property of element matrix.
for i = 1:1:basis_Degree + 1
    for j = 1:1:i - 1
        K_Elem(i, j) = K_Elem(j, i);
    end;
end;

\end{lstlisting}
